\documentclass[10pt,a4paper]{amsart}
\usepackage[margin=19mm]{geometry}
\usepackage{amsmath,amssymb,mathtools}
\usepackage[scr=rsfs,cal=euler]{mathalfa}
\usepackage{xfrac}
\usepackage[unicode,hidelinks,bookmarks=false]{hyperref}
\newcommand{\ie}{i.e.\ }
\newcommand{\cf}{cf.\ }
\newcommand{\resp}{respectively\ }
\newcommand{\Subsection}{\subsection}
\providecommand{\nobreakdash}{\nobreak\hbox{-}}
\setlength{\emergencystretch}{2em}
\multlinegap0pt
\usepackage{amssymb}
\usepackage[shortlabels]{enumitem}
\newtheorem{prop}{Proposition}
\def\End{\ensuremath{\operatorname{End}}}
\title{Descent Theory for Vertex Algebras}
\author{Robin Mader, Terry Gannon,, Arturo Pianzola}
\date{Source: arXiv:2511.02251v2 (CC BY 4.0)}
\begin{document}
\maketitle

\noindent\emph{Source and licence.} Descent Theory for Vertex Algebras. Version arXiv:2511.02251v2 (CC BY 4.0), distributed by arXiv under the Creative Commons Attribution 4.0 International licence, http://creativecommons.org/licenses/by/4.0/, as recorded in the arXiv licence metadata for this exact version and shown on the abstract page https://arxiv.org/abs/2511.02251v2. Full licence notice: \texttt{licenses/arxiv-2511.02251-CC-BY-4.0.txt}.\par\smallskip

\noindent\textit{Editorial scope.} Counted unit: the article's prop tfiso (source0.tex lines 706--713). The article's complete original proof of it is delivered with it (source0.tex lines 715--719, one \texttt{proof} environment of 5 lines). No mathematics is written, repaired or re-derived: the statement and the proof are the article's own lines, in the article's own notation, and no retained line is substituted or re-flowed.  Retained uncounted as the source-local context the proof needs: lines 700--700.  Nothing was omitted from any retained span, so the package carries no omission record.  No retained line contains a citation, so the wrapper carries no bibliography and no bibliography entry.  No retained line contains a figure environment, an image-inclusion command or drawing code, so no image is delivered and the package declares no asset dependency.  Editorial formatting only, in addition: an AMS \texttt{amsart} preamble that restates the article's own declarations and macros used by the retained lines, namely the environments \texttt{prop} on separate counters, together with the standard packages those lines need.  The retained environments print in this document as Proposition 1 in order of appearance, whereas the article prints the same items with the numbering of its own full document.  The complete native source of the article as distributed in the arXiv e-print for this exact version is bundled unchanged as \texttt{sources/188394/source0.tex}.
\begin{equation} \label{keyf} (u_{(m)})_\ell(v_{(n)}) = \sum_{i \geq 0} u_{\ell+i} v \otimes (t^m)_{-i-1} t^n = \sum_{i \geq 0} \binom{m}{i} (u_{\ell+i} v)_{(m+n-i)} \end{equation}

\begin{prop}\label{tfiso}
    Let $h \colon V \to V$ be another automorphism of finite period $M$. Let $\varphi \colon L(V,g) \to L(V,h)$ be a $k[t^{\pm1}]-$vertex algebra map, and let $W$ be a weak $h-$twisted $V-$module with vertex operator map $Y_W(v,z) = \sum_{n \in \frac{1}{M} \mathbb Z} v_n z^{-n-1}$. 
    Denote by $\overline \varphi$ the linear map $L(V,g) \to \End(W)$ obtained by composing $\varphi$ with $L(V,h) \to \End(W)$, $v_{(n)} \mapsto v_n$. 
    
    Then the linear map $(\varphi^\ast Y_W)(\cdot, z) \colon V \to \End(W)[[z^{\frac{1}{M}},z^{-\frac{1}{M}}]]$ defined by 
	   \[ (\varphi^\ast Y_W) (v,z) = \sum_{n \in \frac{r}{M} + \mathbb Z} \overline \varphi(v_{(n)}) z^{-n-1}, \quad v \in V^{g,r}, \]
    makes $W$ into a weak $g-$twisted $V-$module, which we denote simply by $\varphi^\ast W$.
\end{prop}

\begin{proof}
	For $u \in V^{g,r}$, $m \in \frac{r}{M} + \mathbb Z$, we write $\varphi(u_{(m)}) = \sum_j u^j_{(m_j)}$ for some $u^j \in V^{h, r_j}$ and $m_j \in \frac{r_j}{M} + \mathbb Z$, and observe then that $\varphi (u_{(m+\ell)}) = \sum_j u^j_{(m_j+\ell)}$ follows for any $\ell \in \mathbb Z$ from $k[t^{\pm 1}]-$linearity.
	The regularity axiom follows, and so does the vacuum axiom, since $\varphi(\mathbf 1_{(0)}\!) = \mathbf 1_{(0)}$.
	Verification of the Jacobi identity  for $\varphi^\ast Y_W$ additionally uses \eqref{keyf}, the twisted Jacobi identity for $Y_W$, and that $\varphi$ is a homomorphism, but is otherwise mechanical.
\end{proof}

\end{document}
